\documentclass[a4paper]{article}

% Basic packages
% Español (México) con XeLaTeX: fontspec para fuentes Unicode, polyglossia
% para la división silábica y los nombres del documento en la variante mexicana.
\usepackage{fontspec}
\usepackage{polyglossia}
\setdefaultlanguage[variant=mexican]{spanish}
% Los nombres chinos van como «pinyin (original)»: xeCJK da a los caracteres
% Han su propia fuente; sin ella XeLaTeX los omite con un aviso.
% FandolSong no cubre algunos caracteres de nombres (U+73FA); WenQuanYi Zen Hei sí.
\usepackage[AutoFallBack=true]{xeCJK}
\setCJKmainfont{FandolSong-Regular.otf}
\setCJKfallbackfamilyfont{\CJKrmdefault}{WenQuanYi Zen Hei}
% polyglossia no traduce los nombres de listings.
\AtBeginDocument{\providecommand{\lstlistingname}{}\renewcommand{\lstlistingname}{Listado}%
  \providecommand{\lstlistlistingname}{}\renewcommand{\lstlistlistingname}{Índice de listados}}
\usepackage{amsmath, amssymb}
\usepackage{graphicx}
\usepackage[margin=2.5cm]{geometry}
\usepackage[most]{tcolorbox}
\usepackage{etoolbox}
\usepackage{listings}
\usepackage{hyperref}
\usepackage{booktabs}
\usepackage{subcaption}
\usepackage{float}

% Highlight boxes
\newtcolorbox{knowledgebox}[1]{
    enhanced, colback=blue!5!white, colframe=blue!75!black, colbacktitle=blue!75!black,
    coltitle=white, fonttitle=\bfseries, title={#1}, attach boxed title to top left={yshift=-2mm, xshift=2mm},
    boxrule=1pt, sharp corners
}

\newtcolorbox{importantbox}[1]{
    enhanced, colback=yellow!10!white, colframe=yellow!80!black, colbacktitle=yellow!80!black,
    coltitle=black, fonttitle=\bfseries, title={#1}, sharp corners
}

\newtcolorbox{warningbox}[1]{
    enhanced, colback=red!5!white, colframe=red!75!black, colbacktitle=red!75!black,
    coltitle=white, fonttitle=\bfseries, title={#1}, sharp corners
}

% Listings
\lstset{
    language=Python,
    basicstyle=\ttfamily\small,
    keywordstyle=\color{blue},
    stringstyle=\color{red!60!black},
    commentstyle=\color{green!60!black},
    breaklines=true,
    frame=single,
    numbers=left,
    numberstyle=\tiny\color{gray},
    captionpos=b,
    extendedchars=true
}

% Front-page metadata
\newcommand{\notetitle}{MIT 6.S191 Lecture 8: Large Language Models --- Post-Training}
\newcommand{\noteauthors}{Elaboradas a partir de materiales públicos del curso}
\newcommand{\notedate}{Primavera de 2025}
\newcommand{\videochannel}{Liquid AI (Guest Lecture)}
\newcommand{\videopublishdate}{2025-04-21}
\newcommand{\videoduration}{68 minutos}
\newcommand{\videourl}{https://www.youtube.com/watch?v=_HfdncCbMOE}
\newcommand{\videocoverpath}{cover.jpg}
\newcommand{\repourl}{https://github.com/hqhq1025/ai-course-notes}

% Footer with repo info
\usepackage{fancyhdr}
\pagestyle{fancy}
\fancyhf{}
\fancyfoot[C]{\thepage}
\fancyfoot[R]{\footnotesize\color{gray}\href{\repourl}{AI Course Notes}}
\renewcommand{\headrulewidth}{0pt}
\renewcommand{\footrulewidth}{0pt}

\begin{document}

\begin{titlepage}
\centering
{\Large Notas del curso\par}
\vspace{1.2cm}
{\huge\bfseries \notetitle\par}
\vspace{0.8cm}
{\large \noteauthors\par}
\vspace{0.3cm}
{\large \notedate\par}
\vspace{1.2cm}

\ifdefempty{\videocoverpath}{
{\small Al generar las notas, indique la ruta local de la portada del video.\par}
}{
\includegraphics[width=0.82\textwidth,height=0.45\textheight,keepaspectratio]{\videocoverpath}\par
}

\vfill
\begin{tcolorbox}[width=0.9\textwidth, colback=black!2!white, colframe=black!60, sharp corners]
\textbf{Autor o canal del video}: \videochannel\par
\textbf{Fecha de publicación}: \videopublishdate\par
\textbf{Duración del video}: \videoduration\par
\ifdefempty{\videourl}{}{
\textbf{Enlace del video}: \href{\videourl}{\nolinkurl{\videourl}}\par
}\par
\textbf{Página del proyecto}: \href{\repourl}{\nolinkurl{\repourl}}\par
\end{tcolorbox}
\end{titlepage}

\tableofcontents
\newpage

%% --- Inicio del contenido --- %%

\section{Introducción: qué es el Post-Training}

Esta clase la imparten investigadores de Liquid AI (autores de «LLM Engineer's Handbook»), quienes presentan de manera sistemática el flujo de Post-Training (posentrenamiento) de los modelos de lenguaje de gran escala (LLM). El Post-Training abarca todos los pasos de entrenamiento posteriores al preentrenamiento, y su objetivo es convertir un base model que solo sabe predecir el siguiente token en un AI assistant realmente útil.

Durante la etapa de preentrenamiento, el modelo se entrena sobre una cantidad masiva de texto original con el objetivo de next-token prediction; el base model que resulta tiene una gran capacidad de modelado del lenguaje, pero carece de la capacidad de responder preguntas o seguir instrucciones. El Post-Training es precisamente el paso clave que cierra esta brecha.

\begin{importantbox}{Las dos etapas centrales del Post-Training}
\begin{enumerate}
    \item \textbf{Supervised Fine-Tuning (SFT)}: mediante pares de instrucción y respuesta (instruction-answer pairs) se le enseña al modelo a seguir instrucciones y a responder preguntas en forma de conversación. En esta etapa el modelo aprende la estructura del chat template (system / user / assistant).
    \item \textbf{Preference Alignment (alineación de preferencias)}: se ajusta el comportamiento del modelo comparando la chosen answer (respuesta deseada) con la rejected answer (respuesta no deseada), de modo que las respuestas producidas se ajusten mejor a las preferencias humanas.
\end{enumerate}
\end{importantbox}

\begin{knowledgebox}{Diferencia de escala de datos entre Pre-Training y Post-Training}
El preentrenamiento utiliza decenas de billones (trillions) de tokens de texto original, con un conjunto de datos de tamaño extremadamente grande. En cambio, el conjunto de datos del Post-Training es mucho más pequeño, por lo general del orden de miles a millones de ejemplos, pero exige una \textbf{calidad} de los datos mucho mayor que la cantidad. Esta diferencia es la base para entender la metodología del Post-Training.
\end{knowledgebox}

\subsection{Resumen de la sección}

El Post-Training es el paso central que convierte un base model en un assistant utilizable, e incluye dos etapas principales: SFT y Preference Alignment. Se caracteriza por tener un conjunto de datos mucho más pequeño que el del preentrenamiento, pero con exigencias de calidad muy altas.

\newpage
\section{Tipos de Fine-Tuning y escenarios de aplicación}

\subsection{Tres tipos de Fine-Tuning}

El ponente divide el Fine-Tuning en tres categorías, cada una con distintas necesidades de escala de datos y objetivos:

\begin{table}[H]
\centering
\caption{Comparación de tipos de Fine-Tuning}
\begin{tabular}{lccc}
\toprule
\textbf{Tipo} & \textbf{Objetivo} & \textbf{Escala de datos} & \textbf{Ejemplo} \\
\midrule
Propósito general (General Purpose) & Crear un chatbot polivalente & $>$1,000,000 muestras & ChatGPT, Claude \\
Específico de dominio (Domain Specific) & Incorporar conocimiento de dominio & miles$\sim$cientos de miles & medicina/derecho/finanzas \\
Específico de tarea (Task Specific) & Aprender una única función & cientos$\sim$miles & resumen, corrección ortográfica \\
\bottomrule
\end{tabular}
\end{table}

\begin{warningbox}{La necesidad de datos depende de múltiples factores}
Las cantidades de muestras anteriores son solo estimaciones aproximadas. La necesidad real depende de: (1) \textbf{el tamaño del modelo}: cuanto más grande es el modelo, menos muestras se necesitan; (2) \textbf{la complejidad de la tarea}: si el conocimiento del base model sobre el idioma o dominio objetivo es cercano a cero (como una lengua poco común que nunca ha visto), el efecto del fine-tuning será muy limitado, porque no hay suficientes tokens para «aprender desde cero».
\end{warningbox}

\subsection{Cuándo conviene hacer Fine-Tuning}

El ponente recomienda seguir un proceso de decisión claro:

\begin{enumerate}
    \item \textbf{Probar primero In-Context Learning / RAG}: es la opción con la barrera más baja y permite una validación rápida.
    \item Si el resultado no es satisfactorio, considerar las siguientes cuatro motivaciones para el fine-tuning:
    \begin{itemize}
        \item \textbf{Cambiar el estilo o el formato de la salida}: por ejemplo, ajustar el tono con el que el modelo genera correos electrónicos
        \item \textbf{Agregar conocimiento de dominio} (superficial): por ejemplo, enseñarle al modelo hechos sobre tu empresa
        \item \textbf{Destilación de conocimiento}: destilar capacidades del nivel de GPT-4 en un modelo más pequeño, para reducir el costo y la latencia
        \item \textbf{Mejorar la calidad en una tarea específica}: por ejemplo, en la generación de diagram, es posible superar a un frontier model en una tarea acotada
    \end{itemize}
\end{enumerate}

\begin{knowledgebox}{Motivaciones comerciales de las empresas para elegir Fine-Tuning}
Además del aspecto técnico, hay dos razones importantes no técnicas por las que las empresas eligen la ruta de modelo de código abierto más fine-tuning: (1) \textbf{control de los datos}: no quieren enviar datos sensibles a la nube mediante una API; (2) \textbf{personalización}: quieren personalizar a fondo el comportamiento, el tono y el formato de salida del modelo, para crear un modelo propio.
\end{knowledgebox}

\subsection{Resumen de la sección}

El Fine-Tuning se divide en tres niveles: propósito general, específico de dominio y específico de tarea, con una necesidad de datos que disminuye sucesivamente. En la práctica, primero se debe probar in-context learning y, solo después de confirmar que es insuficiente, considerar el fine-tuning. El valor del Fine-Tuning no está solo en la mejora técnica, sino también en dar a las empresas el control sobre el modelo y los datos.

\newpage
\section{Conjuntos de datos de post-entrenamiento: tres dimensiones de la calidad}

\subsection{Definición de buenos datos}

Quien imparte la clase considera que la calidad de los datos es el problema más central de todo el proceso de post-entrenamiento, y puede medirse a partir de tres dimensiones:

\begin{importantbox}{Marco tridimensional de la calidad de los datos}
\begin{enumerate}
    \item \textbf{Accuracy (exactitud)}: si la respuesta es correcta. Es relativamente fácil de verificar para preguntas simples; para código puede verificarse mediante unit tests, y para matemáticas mediante un solver.
    \item \textbf{Diversity (diversidad)}: si el conjunto de ejemplos cubre los distintos escenarios de interacción que puede tener el usuario. Una diversidad insuficiente provoca que al modelo le falte gravemente capacidad en ciertos dominios.
    \item \textbf{Complexity (complejidad)}: si los ejemplos son lo bastante desafiantes, con razonamiento de varios pasos, chain-of-thought, etc.
\end{enumerate}
\end{importantbox}

\subsection{La importancia de la diversidad}

Al visualizar en 2D el embedding de distintos conjuntos de datos, la diferencia se aprecia de manera directa:

\begin{itemize}
    \item \textbf{Conversaciones reales de usuarios} (como ShareGPT, WildChat): la distribución del embedding cubre un rango sumamente amplio, porque las preguntas de los usuarios abarcan todo tipo de temas
    \item \textbf{Conjuntos de datos matemáticos}: el embedding se concentra en gran medida en una sola región
    \item \textbf{Conjuntos de datos de código}: también se concentran en una región limitada
\end{itemize}

Por lo tanto, si el objetivo es construir un LLM de propósito general, los datos de matemáticas y de código por sí solos son claramente insuficientes; es necesario incorporar datos de conversación diversos para ampliar la cobertura.

\begin{warningbox}{La trampa de la diversidad en los datos sintéticos}
Uno de los grandes riesgos de los datos sintéticos (synthetic data) es el colapso de la diversidad. Una pequeña cantidad de datos sintéticos puede ampliar la diversidad de manera efectiva, pero un \textbf{exceso} de datos sintéticos, por el contrario, provoca que la distribución de los datos se estreche y que la diversidad caiga de forma abrupta, lo que termina por reducir la calidad general de los datos.
\end{warningbox}

\subsection{Aumento de la complejidad: Auto-Evol}

Un ejemplo de baja complejidad como «¿Qué tan alta es la Torre Eiffel?» puede responderse con solo unas cuantas palabras. Es posible usar marcos como \textbf{Auto-Evol} para que el LLM reescriba el prompt y aumente su complejidad; por ejemplo, exigiendo un análisis desde varios ángulos, comparando datos de distintas épocas, explicando principios de ingeniería, etc. Las preguntas reescritas requieren respuestas más largas y detalladas, lo que aumenta de manera natural la complejidad de los datos.

\subsection{Formato de los datos}

El post-entrenamiento utiliza dos formatos de datos principales:

\begin{table}[H]
\centering
\caption{Formatos de datos de post-entrenamiento}
\begin{tabular}{lll}
\toprule
\textbf{Formato} & \textbf{Composición} & \textbf{Uso} \\
\midrule
Instruction Data & system prompt + instruction + output & SFT \\
Preference Data & instruction + chosen answer + rejected answer & DPO/PPO \\
\bottomrule
\end{tabular}
\end{table}

\subsection{Resumen de la sección}

La calidad de los datos está determinada en conjunto por tres dimensiones: exactitud, diversidad y complejidad. Al construir un conjunto de datos, hay que prestar especial atención al riesgo de diversidad de los datos sintéticos, y pueden usarse herramientas como Auto-Evol para aumentar la complejidad de los ejemplos.

\newpage
\section{Pipeline de generación de conjuntos de datos}

\subsection{Generación de Instruction Data}

El expositor presentó un pipeline típico de generación de instruction data:

\begin{enumerate}
    \item \textbf{Seed Data}: el punto de partida puede ser texto original, pares de preguntas y respuestas (QA) ya existentes, o consultas reales de usuarios
    \item \textbf{Generación de instrucciones}: si solo se cuenta con texto original (es decir, la respuesta), se puede usar \textbf{back translation} (traducción inversa) para que el LLM genere la pregunta correspondiente
    \item \textbf{Generación de respuestas}: se usa un LLM para generar la respuesta a partir de la instrucción
    \item \textbf{Evaluación y filtrado}:
    \begin{itemize}
        \item Se evalúa la calidad con reglas heurísticas o con LLM-as-Judge
        \item Deduplicación de los datos (deduplication)
        \item Descontaminación (decontamination): garantizar que no se incluyan muestras del conjunto de prueba
        \item Filtrado por palabras clave: retirar contenido no deseado
    \end{itemize}
\end{enumerate}

\subsubsection{Tratamiento especial de los datos de Instruction Following}

Instruction following es un tipo de datos centrado en el \textbf{cumplimiento de restricciones}. Por ejemplo: «responde en inglés y todo en minúsculas». Al generar este tipo de datos, se puede anexar la condición de restricción al final del prompt para que el LLM genere la respuesta, y luego usar \textbf{pruebas programáticas} para verificar si se cumple la restricción (como detección de idioma o verificación de mayúsculas y minúsculas), descartando las muestras que no cumplan los requisitos.

\subsection{Generación de Preference Data: el paradigma Ultra Feedback}

La generación de Preference data sigue una estrategia distinta:

\begin{enumerate}
    \item Dado un prompt, se consulta a \textbf{varios LLM distintos} (en lugar de un solo modelo)
    \item Se usa un Judge LLM para calificar cada respuesta
    \item Se toma la respuesta con la calificación más alta como \textbf{chosen} y la de calificación más baja como \textbf{rejected}
    \item Deduplicación, retiro de respuestas demasiado cortas y otros procesamientos posteriores
\end{enumerate}

\begin{knowledgebox}{Decontamination (descontaminación)}
La descontaminación es un paso fundamental en el pipeline de generación de datos. Por ejemplo, IFEval es un conjunto de evaluación de instruction following ampliamente usado. Al generar datos de entrenamiento, es indispensable garantizar que las muestras generadas no tengan una alta similitud con el conjunto de prueba; de lo contrario, la puntuación de evaluación se inflará de forma artificial sin que la capacidad real mejore.
\end{knowledgebox}

\subsection{Caso: el conjunto de datos Open Perfect Blend}

El expositor mostró el conjunto de datos Open Perfect Blend que elaboró, un conjunto de datos mixto orientado al fine-tuning general. Su proporción aproximada por categoría es:

\begin{itemize}
    \item \textbf{Math} (matemáticas): proporción mayor
    \item \textbf{Code} (código): proporción mayor
    \item \textbf{Chat} (diálogo general): proporción media
    \item \textbf{Instruction Following}: proporción menor
\end{itemize}

El diseño de esta proporción de mezcla (data mixture) determina directamente el desempeño del modelo en las distintas dimensiones de capacidad.

\subsection{Chat Template}

Una vez terminada la preparación de los datos, el último paso es aplicar el \textbf{Chat Template}. Esta es la estructura clave que permite al modelo entender «quién está hablando»:

\begin{lstlisting}[caption={Ejemplo de Chat Template (formato ChatML)}]
<|im_start|>system
You are a helpful assistant.
<|im_end|>
<|im_start|>user
What is the capital of France?
<|im_end|>
<|im_start|>assistant
The capital of France is Paris.
<|im_end|>
\end{lstlisting}

Los tokens especiales \texttt{<|im\_start|>} y \texttt{<|im\_end|>} marcan el inicio y el fin de cada mensaje, y las etiquetas de rol (system / user / assistant) precisan la identidad de quien habla. Las conversaciones de varios turnos se logran repitiendo el par user-assistant.

\begin{importantbox}{El Chat Template es el objetivo central de aprendizaje del SFT}
La estructura del Chat Template es uno de los contenidos centrales que el modelo debe aprender en la etapa de SFT. Es precisamente este formato de conversación estructurado el que hace que el base model evolucione de «solo saber continuar texto» a «un assistant capaz de sostener conversaciones de varios turnos». Distintas familias de modelos (Llama, Mistral, Qwen, etc.) usan chat templates diferentes, pero el principio es el mismo.
\end{importantbox}

\subsection{Resumen de la sección}

La generación de conjuntos de datos es un pipeline sistemático que incluye la preparación de los datos semilla, la generación de instrucciones y respuestas, la evaluación y el filtrado, y la deduplicación y descontaminación, entre otros pasos. Los datos de preferencia (Preference data) se obtienen mediante muestreo con varios modelos más la calificación de un Judge. Finalmente, los datos deben formatearse con el Chat Template para usarse en el entrenamiento de SFT.

\newpage
\section{Algoritmos y técnicas de entrenamiento}

\subsection{Bibliotecas recomendadas de fine-tuning}

El presentador recomendó tres bibliotecas principales de fine-tuning:

\begin{table}[H]
\centering
\caption{Comparación de las principales bibliotecas de fine-tuning}
\begin{tabular}{llp{7cm}}
\toprule
\textbf{Herramienta} & \textbf{Desarrollador} & \textbf{Características} \\
\midrule
TRL & Hugging Face & Basada en Transformers, con la mayor variedad de algoritmos y actualizada con la investigación más reciente \\
Axolotl & Comunidad & Basada en TRL, impulsada por configuración YAML, fácil de compartir y reproducir configuraciones de entrenamiento \\
Unsloth & Comunidad & Optimizada para una sola GPU, con una eficiencia muy alta e incorpora funciones prácticas como la cuantización \\
\bottomrule
\end{tabular}
\end{table}

\subsection{Las tres técnicas de entrenamiento de SFT}

\begin{importantbox}{Full Fine-Tuning vs. LoRA vs. QLoRA}
\begin{itemize}
    \item \textbf{Full Fine-Tuning}: carga el modelo con precisión completa y entrena el 100\% de los parámetros. Ofrece la mayor calidad, pero requiere una cantidad enorme de VRAM y, por lo general, un clúster de GPU completo.
    \item \textbf{LoRA (Low-Rank Adaptation)}: congela los pesos originales del modelo y solo entrena las matrices adicionales de adaptador de bajo rango A y B. Solo necesita entrenar alrededor del 0.5\% de los parámetros, es rápido y requiere poca VRAM; \textbf{se recomienda como opción predeterminada}.
    \item \textbf{QLoRA (Quantized LoRA)}: primero cuantiza el modelo a una precisión de 4 bits y luego lo carga, para después aplicar LoRA. Reduce aún más los requisitos de VRAM, pero el entrenamiento es más lento y el rendimiento presenta una ligera caída (unos cuantos puntos porcentuales).
\end{itemize}
\end{importantbox}

Principio matemático de LoRA: para la matriz de pesos original $W \in \mathbb{R}^{d \times k}$, LoRA agrega una descomposición de bajo rango:

$$W' = W + \Delta W = W + BA$$

donde:
\begin{itemize}
    \item $B \in \mathbb{R}^{d \times r}$, $A \in \mathbb{R}^{r \times k}$
    \item $r \ll \min(d, k)$ es el rango (rank), y normalmente se elige entre 8 y 64
    \item Durante el entrenamiento solo se actualizan $A$ y $B$; el $W$ original permanece congelado
\end{itemize}

\begin{warningbox}{Full Fine-Tuning suele ser excesivo}
A menos que seas una empresa dedicada específicamente a entrenar LLM, LoRA es casi siempre la mejor opción. Full Fine-Tuning puede provocar sobreajuste, y la mejora en el rendimiento a menudo no justifica el costo computacional adicional. Solo si los requisitos de VRAM de LoRA siguen siendo demasiado altos conviene recurrir a QLoRA.
\end{warningbox}

\subsection{Algoritmos de Preference Alignment}

El campo de la alineación de preferencias cuenta con más de 100 algoritmos distintos, pero el presentador recomienda memorizar solo dos:

\subsubsection{PPO (Proximal Policy Optimization)}

PPO es el algoritmo clásico de RLHF y requiere cargar \textbf{tres modelos} simultáneamente:

\begin{enumerate}
    \item \textbf{Frozen Model} (política de referencia): una instantánea del modelo previa al entrenamiento, que se usa para calcular la KL divergence
    \item \textbf{Train Model} (política actual): el modelo que recibe el prompt y genera la respuesta
    \item \textbf{Reward Model}: califica la respuesta generada
\end{enumerate}

Flujo de entrenamiento: el Train Model genera una respuesta $\rightarrow$ el Reward Model la califica $\rightarrow$ se usa la KL divergence para restringir la magnitud de la actualización de la política $\rightarrow$ se actualizan los parámetros del Train Model.

La ventaja es que el límite superior de calidad es alto; la desventaja es que requiere tres modelos, es extremadamente costoso y difícil de ajustar.

\subsubsection{DPO (Direct Preference Optimization)}

DPO solo necesita \textbf{dos modelos} (si se usa LoRA, en realidad basta con cargar un modelo base más dos conjuntos de adaptadores):

$$\mathcal{L}_{\text{DPO}} = -\mathbb{E}\left[\log \sigma\left(\beta \log \frac{\pi_\theta(y_w|x)}{\pi_{\text{ref}}(y_w|x)} - \beta \log \frac{\pi_\theta(y_l|x)}{\pi_{\text{ref}}(y_l|x)}\right)\right]$$

donde:
\begin{itemize}
    \item $y_w$ es la chosen answer y $y_l$ es la rejected answer
    \item $\pi_\theta$ es la política en entrenamiento y $\pi_{\text{ref}}$ es la política de referencia
    \item $\beta$ controla la intensidad de la penalización por desviarse de la política de referencia
\end{itemize}

\begin{importantbox}{Recomendación práctica: DPO es mejor que PPO}
A menos que se cuente con recursos prácticamente ilimitados al nivel de OpenAI, se recomienda usar DPO. DPO es más rápido, más económico y más fácil de ajustar, con una calidad apenas ligeramente inferior a la de PPO. DPO no necesita un reward model independiente, sino que optimiza directamente sobre preference data, lo cual simplifica enormemente el flujo de entrenamiento.
\end{importantbox}

\subsection{Hiperparámetros clave del entrenamiento}

\begin{table}[H]
\centering
\caption{Hiperparámetros clave del Post-Training}
\begin{tabular}{lll}
\toprule
\textbf{Parámetro} & \textbf{Valor recomendado} & \textbf{Descripción} \\
\midrule
Learning Rate & $1 \times 10^{-5}$ $\sim$ $5 \times 10^{-5}$ & El parámetro más crítico; si es demasiado alto provoca un loss spike \\
Batch Size & Limitado por la VRAM & Se puede usar gradient accumulation para simular un batch grande \\
Max Length & Limitado por la VRAM & Determina la longitud máxima de la secuencia de entrada \\
Epochs & 3$\sim$5 & Número de veces que se recorre por completo el conjunto de entrenamiento \\
Optimizer & AdamW & Una base sólida que normalmente no hace falta cambiar \\
Attention & Flash Attention & Una implementación eficiente de la atención que acelera especialmente el procesamiento de secuencias largas \\
\bottomrule
\end{tabular}
\end{table}

\begin{knowledgebox}{Gradient Accumulation y Effective Batch Size}
Cuando la memoria de la GPU no basta para admitir el batch size deseado, se puede usar gradient accumulation: se acumulan los gradientes tras varios forward pass y luego se actualizan los pesos de una sola vez. De este modo, \textbf{effective batch size} = per-GPU batch size $\times$ accumulation steps $\times$ número de GPU, lo que permite simular un batch más grande sin aumentar la memoria.
\end{knowledgebox}

\subsection{Monitoreo del entrenamiento}

Durante el entrenamiento es necesario monitorear tres métricas fundamentales:

\begin{enumerate}
    \item \textbf{Training Loss}: debe disminuir de forma uniforme; si aparece un loss spike, por lo general significa que el learning rate es demasiado alto
    \item \textbf{Validation Loss}: se usa entre el 1 y el 2\% de los datos como conjunto de validación para monitorear si hay sobreajuste
    \item \textbf{Gradient Norm}: la magnitud del gradiente; demasiados spike pueden indicar un problema con los datos
\end{enumerate}

\begin{warningbox}{Diagnóstico de un Loss Spike}
Si al principio del entrenamiento el loss disminuye de forma estable, pero luego aparece de pronto un salto brusco (spike), la causa más probable es que el \textbf{learning rate sea demasiado alto}. La solución es reducir el learning rate, lo cual debería producir una curva de loss más uniforme. Ten cuidado de no confundir las pequeñas fluctuaciones del inicio del entrenamiento con un spike — las fluctuaciones iniciales de poca magnitud suelen ser normales.
\end{warningbox}

\subsection{Resumen de la sección}

Para SFT se recomienda usar LoRA como técnica de entrenamiento predeterminada, y para Preference Alignment se recomienda usar DPO. Entre los hiperparámetros clave, el learning rate es el más importante, y AdamW junto con Flash Attention es una opción predeterminada confiable. Durante el entrenamiento se debe monitorear de cerca el training loss, el validation loss y el gradient norm.

\newpage
\section{Model Merging: mejora de modelos sin entrenamiento}

\subsection{Conceptos básicos}

Model Merging es una técnica que combina directamente los parámetros de varios modelos, \textbf{sin necesidad de entrenamiento adicional}. Al principio se consideraba una «broma», pero ahora la ha adoptado toda empresa de LLM.

Idea central: tomar varios modelos de la misma arquitectura que pasaron por distintos fine-tuning y combinar sus parámetros mediante operaciones matemáticas (como el promedio ponderado), para obtener un nuevo modelo que reúna las ventajas de todos.

\begin{warningbox}{Limitaciones del Model Merging}
\begin{itemize}
    \item Solo pueden combinarse modelos de \textbf{la misma arquitectura y el mismo tamaño} (por ejemplo, solo se puede combinar Llama 3 con Llama 3, no entre generaciones distintas)
    \item Se recomienda usar modelos de \textbf{precisión completa} para la combinación: la pérdida de precisión de los modelos cuantizados se propaga al resultado combinado
    \item Elegir los pesos y parámetros adecuados se parece más a la «alquimia» que a una ciencia rigurosa: requiere experimentación repetida
\end{itemize}
\end{warningbox}

\subsection{Principales técnicas de combinación}

\subsubsection{SLERP (Spherical Linear Interpolation)}

SLERP realiza la interpolación sobre una esfera (en lugar de una interpolación lineal), lo que normalmente produce mejores resultados de combinación que el simple promedio lineal.

$$\text{SLERP}(\mathbf{p}_0, \mathbf{p}_1; t) = \frac{\sin((1-t)\Omega)}{\sin\Omega}\mathbf{p}_0 + \frac{\sin(t\Omega)}{\sin\Omega}\mathbf{p}_1$$

donde $\Omega = \arccos(\mathbf{p}_0 \cdot \mathbf{p}_1 / \|\mathbf{p}_0\| \|\mathbf{p}_1\|)$, $t \in {[}0, 1{]}$ es el coeficiente de interpolación.

Características: sencillo y confiable, muy popular en la comunidad de código abierto; pero \textbf{solo puede combinar dos modelos}.

\subsubsection{DARE + TIES}

Una técnica de combinación más avanzada, capaz de procesar más de dos modelos:

\begin{itemize}
    \item \textbf{DARE (Drop And REscale)}: poda aleatoriamente y reescala los parámetros de los modelos de origen
    \item \textbf{TIES (TrIm, Elect Sign, merge)}: conserva los parámetros más significativos y resuelve el problema de la cancelación de signos entre parámetros positivos y negativos mediante el \textbf{sign consensus} (consenso de signos)
\end{itemize}

Motivación del sign consensus: si un parámetro del modelo A es positivo y el mismo parámetro del modelo B es negativo, el simple promedio da un resultado cercano a cero, lo que en realidad hace perder la información que cada modelo aprendió por separado. TIES determina el signo del parámetro mediante un mecanismo de votación, evitando esta cancelación.

\subsection{Configuración práctica de Model Merging}

Con herramientas como \textbf{mergekit}, el esquema de combinación se define mediante un archivo de configuración YAML:

\begin{itemize}
    \item \textbf{density}: controla la proporción de parámetros que se conservan (en el método TIES no se conservan todos los parámetros)
    \item \textbf{weight}: controla el peso de cada modelo de origen en la combinación (se puede asignar un peso distinto según la puntuación de evaluación)
\end{itemize}

El expositor mostró el árbol genealógico del modelo Daredevil 8B: mediante combinación iterativa (un modelo combinado se vuelve a combinar con otros modelos) se puede construir un linaje de modelos complejo, y el producto final supera a todos los modelos de origen en varios benchmarks.

\subsection{Caso de aplicación: modelos específicos de idioma}

Proceso completo tomando como ejemplo la construcción de un LLM en finés:

\begin{enumerate}
    \item Preentrenamiento continuo sobre texto original en finés (5-100B tokens)
    \item SFT sobre instruction data en finés
    \item Preference Alignment sobre preference data en finés
    \item Problema: la capacidad del modelo en finés mejora mucho, pero \textbf{todas las demás capacidades se degradan gravemente}
    \item Solución: combinar el modelo en finés con un modelo instruct general (como Llama 3 Instruct)
    \item Resultado: se obtiene un modelo que \textbf{domina el finés y conserva las capacidades generales}
\end{enumerate}

\begin{importantbox}{El valor central de Model Merging}
Model Merging permite \textbf{superponer} las especialidades de distintos modelos sin sacrificar las capacidades que cada uno ya tenía. Es una estrategia eficaz para resolver el problema del «olvido catastrófico» (catastrophic forgetting) durante el fine-tuning, y resulta especialmente adecuada para escenarios como modelos específicos de idioma o modelos específicos de dominio.
\end{importantbox}

\subsection{Resumen de la sección}

Model Merging es una técnica de mejora de modelos sin costo de entrenamiento, que combina las ventajas de varios modelos mediante operaciones matemáticas a nivel de parámetros. SLERP es adecuado para combinar dos modelos, mientras que DARE+TIES admite la combinación de varios modelos. Es una herramienta poderosa para resolver el problema de degradación de capacidades causado por el fine-tuning.

\newpage
\section{Evaluación del modelo}

\subsection{La importancia de la evaluación}

El expositor admite que la evaluación de los LLM es uno de los mayores problemas de todo el campo: «en realidad no sabemos muy bien qué estamos haciendo». Pero la evaluación sigue siendo fundamental, porque la esencia del Post-Training es la \textbf{optimización}, y si la métrica de evaluación no es la correcta, la dirección de la optimización se desvía.

\subsection{Método de evaluación uno: Automated Benchmarks}

\textbf{Principio}: dado un conjunto de preguntas (como las de opción múltiple de MMLU), el modelo elige una respuesta y se le da una calificación con métricas como la exactitud.

\begin{table}[H]
\centering
\caption{Benchmarks automatizados comunes}
\begin{tabular}{lll}
\toprule
\textbf{Benchmark} & \textbf{Dominio} & \textbf{Métrica} \\
\midrule
MMLU & Conocimiento general (57 disciplinas) & Exactitud \\
GSM8K & Razonamiento matemático de primaria & Exactitud \\
MATH & Matemáticas de competencia & Exactitud \\
HumanEval & Generación de código & pass@k \\
IFEval & Instruction Following & Tasa de cumplimiento de restricciones \\
\bottomrule
\end{tabular}
\end{table}

\textbf{Ventajas}: escalable, de bajo costo, reproducible, se puede diseñar para tareas específicas.

\textbf{Desventajas}: no refleja el uso real del modelo; el usuario \textbf{conversa} con el modelo, no responde preguntas de opción múltiple.

\subsection{Método de evaluación dos: Human Evaluation (evaluación humana)}

\textbf{Principio}: evaluadores humanos eligen la mejor respuesta entre las de dos modelos anónimos (como en Chatbot Arena).

\textbf{Ventajas}:
\begin{itemize}
    \item Refleja directamente la preferencia humana, que es el objetivo último de optimización del Post-Training
    \item Permite controlar con precisión la dimensión evaluada (por ejemplo, centrarse solo en la toxicidad o solo en la exactitud)
    \item Riesgo bajo de contaminación de datos
\end{itemize}

\textbf{Desventajas}:
\begin{itemize}
    \item Es sumamente difícil de escalar, con costo alto y mucho tiempo
    \item Los evaluadores humanos \textbf{tienen sesgos importantes}: prefieren respuestas más largas y un tono más seguro, aunque el contenido sea incorrecto
\end{itemize}

\begin{warningbox}{La preferencia humana no se correlaciona con los benchmarks automatizados}
El expositor mostró un hallazgo clave: la clasificación de Chatbot Arena tiene una correlación muy baja con benchmarks automatizados como MMLU, GSM8K y HumanEval. Un modelo con puntuación alta en MMLU puede tener un desempeño muy pobre en la preferencia humana, y viceversa. Por eso, \textbf{es necesario usar ambos tipos de evaluación} a la vez para entender de manera integral la capacidad del modelo.
\end{warningbox}

\subsection{Método de evaluación tres: LLM-as-Judge}

Para encontrar un punto medio entre la escalabilidad y la preferencia humana, se puede usar un LLM en lugar de una persona para evaluar:

\begin{itemize}
    \item \textbf{Un solo Judge}: un LLM potente (como GPT-4) califica las respuestas
    \item \textbf{Varios Judge (Jury)}: un jurado formado por varios LLM vota para decidir cuál es mejor, lo que resulta más robusto y confiable
\end{itemize}

\textbf{Ventajas}: se correlaciona mucho con la preferencia humana, puede manejar tareas complejas de formato libre y da retroalimentación directa.

\textbf{Desventajas}: el LLM judge también tiene sesgos propios (y son muy parecidos a los sesgos humanos), por lo que aún se necesita una pequeña cantidad de evaluación humana para verificar la consistencia del judge.

\subsection{Recomendaciones para construir un sistema de evaluación propio}

\begin{enumerate}
    \item \textbf{Empezar lo antes posible}: diseñar la evaluación antes del fine-tuning, de forma similar al desarrollo guiado por pruebas
    \item \textbf{Iterar de manera continua}: el conjunto de datos de evaluación no es algo fijo, hay que corregirlo constantemente según el desempeño del modelo
    \item \textbf{Usarlos de manera combinada}: benchmark automatizado + LLM-as-Judge + una pequeña cantidad de evaluación humana
    \item \textbf{Comparar de forma transversal}: no solo comparar las distintas versiones del propio fine-tune, sino también contrastar con otros modelos y arquitecturas
\end{enumerate}

\begin{knowledgebox}{Desarrollo iterativo guiado por la evaluación}
El expositor enfatiza que la distribución del tiempo en el Post-Training debería ser, aproximadamente: un tercio en la preparación de datos, un tercio en el entrenamiento y un tercio en la evaluación. La evaluación no es una verificación posterior, sino el elemento central que atraviesa todo el proceso. El resultado de cada ronda de evaluación guía directamente la mejora de los datos y el entrenamiento de la siguiente ronda.
\end{knowledgebox}

\subsection{Resumen de la sección}

Existen principalmente tres métodos de evaluación de LLM: el benchmark automatizado (escalable, pero no refleja el uso real), la evaluación humana (la más cercana al objetivo, pero difícil de escalar) y LLM-as-Judge (una solución intermedia). Los tres se complementan y deben usarse de manera combinada. La baja correlación entre la preferencia humana y los benchmarks automatizados es un hecho clave que hay que reconocer al diseñar un sistema de evaluación.

\newpage
\section{Escalado del cómputo en el momento de la inferencia (Test-Time Compute Scaling)}

\subsection{Idea central}

Ya se ha demostrado que el training-time scaling (aumentar los datos de entrenamiento y el cómputo) es eficaz. Test-Time Compute Scaling aplica la misma idea a la etapa de inferencia: \textbf{invertir más cómputo en el momento de la inferencia a cambio de una salida de mayor calidad}.

\subsection{Majority Voting (votación mayoritaria)}

El método más simple de test-time scaling:

\begin{enumerate}
    \item Para la misma pregunta (por ejemplo, un problema de matemáticas), se le pide al modelo que genere $N$ respuestas distintas (usando una temperatura distinta de cero para obtener diversidad)
    \item Se toma la respuesta que aparece con mayor frecuencia como salida final
\end{enumerate}

Intuición: debido a la aleatoriedad que introduce el muestreo top-p, el modelo ocasionalmente elige un token equivocado que hace que toda la cadena de razonamiento se desvíe. Al muestrear varias veces, es muy probable que la respuesta correcta aparezca con mayor frecuencia.

\subsection{Best-of-N (selección óptima)}

Una versión mejorada de Majority Voting:

\begin{enumerate}
    \item Generar $N$ respuestas distintas
    \item Usar un \textbf{Reward Model} o un \textbf{Judge LLM} para puntuar cada respuesta
    \item Elegir la respuesta con la puntuación más alta
\end{enumerate}

En comparación con Majority Voting, Best-of-N aprovecha una señal de evaluación externa, por lo que suele funcionar mejor, aunque requiere ejecutar un modelo de evaluación adicional.

\subsection{Process Reward Model (PRM)}

Un método más refinado: en lugar de evaluar la respuesta completa, evalúa \textbf{cada paso del razonamiento}:

\begin{importantbox}{Flujo de trabajo del Process Reward Model}
\begin{enumerate}
    \item El LLM genera una solución parcial (partial solution)
    \item El PRM puntúa cada paso: la puntuación representa "la probabilidad de que ese paso finalmente conduzca a la respuesta correcta"
    \item Se elige el paso con la puntuación más alta para seguir expandiéndolo
    \item Se repite el proceso anterior para construir la solución completa paso a paso
\end{enumerate}
Esto da lugar a un mecanismo de \textbf{búsqueda en árbol} (tree search): se busca en el espacio de los pasos de razonamiento, y el PRM guía la dirección de la búsqueda.
\end{importantbox}

El profesor demostró este proceso con un ejemplo de multiplicación en base 8:

\begin{itemize}
    \item Attempt 1: el primer paso es correcto, pero el segundo se calcula en base decimal en lugar de en base 8, y el PRM lo rechaza
    \item Attempt 2: tiene un paso correcto más, pero al final comete el mismo error, y el PRM lo rechaza de nuevo
    \item Attempt 4: todos los pasos son correctos, y el PRM lo aprueba
\end{itemize}

\subsection{Resultados experimentales de Scaling}

Los experimentos del equipo de Hugging Face muestran que:

\begin{itemize}
    \item Usando Llama 3.2 1B y 3B (modelos pequeños), al aumentar el número de generaciones por pregunta
    \item en un benchmark de matemáticas, un modelo pequeño + test-time scaling puede \textbf{superar} a Llama 3.1 8B e incluso a Llama 3.1 70B
\end{itemize}

\begin{importantbox}{El equilibrio entre la velocidad de inferencia y la calidad de la salida}
En esencia, Test-Time Compute Scaling intercambia \textbf{tiempo de inferencia} (inference speed) por \textbf{calidad de la salida} (output quality). Esto significa que un modelo pequeño y eficiente junto con un presupuesto de inferencia suficiente puede igualar o incluso superar a modelos de un tamaño mucho mayor en tareas específicas. Esto tiene una importancia considerable para los escenarios con recursos limitados.
\end{importantbox}

\begin{knowledgebox}{De dónde viene el Reward Model}
El Reward Model que se usa en Test-Time Compute Scaling puede reutilizar el reward model entrenado en la etapa de Preference Alignment (dentro del proceso de entrenamiento con PPO). La entrada de este modelo es texto, y su salida es una puntuación de calidad entre 0 y 1. También es posible sustituir el reward model entrenado específicamente por un Judge LLM (mediante prompting).
\end{knowledgebox}

\subsection{Resumen de la sección}

Test-Time Compute Scaling es una de las tendencias más importantes de principios de 2025; su idea central es invertir más cómputo en el momento de la inferencia para mejorar la calidad de la salida. Desde el simple Majority Voting, pasando por Best-of-N, hasta la búsqueda a nivel de paso basada en el Process Reward Model, la complejidad del método y su efectividad aumentan progresivamente. Los experimentos demuestran que un modelo pequeño + test-time scaling puede superar a modelos más grandes.

\newpage
\section{Síntesis y ampliación}

\subsection{Panorama general del proceso de Post-Training}

Todo el Post-Training es un ciclo iterativo, formado por tres etapas de igual peso:

\begin{enumerate}
    \item \textbf{Construcción de datos} (aproximadamente 1/3 del tiempo): diseñar la proporción de mezcla de datos, garantizar exactitud, diversidad y complejidad, y producir en lotes datos de entrenamiento de alta calidad mediante un pipeline de seed data + generación con LLM + filtrado
    \item \textbf{Entrenamiento del modelo} (aproximadamente 1/3 del tiempo): SFT (se recomienda LoRA) $\rightarrow$ Preference Alignment (se recomienda DPO) $\rightarrow$ Model Merging opcional
    \item \textbf{Iteración de evaluación} (aproximadamente 1/3 del tiempo): benchmarks automatizados, LLM-as-Judge y evaluación humana en conjunto, con los resultados de evaluación retroalimentando la siguiente ronda de datos y de mejoras en el entrenamiento
\end{enumerate}

\subsection{Repaso de los puntos centrales}

\begin{importantbox}{Takeaways centrales de esta clase}
\begin{itemize}
    \item El Post-Training convierte un base model en un assistant útil; los pasos centrales son SFT + Preference Alignment
    \item La calidad de los datos importa más que la cantidad de datos, en tres dimensiones: Accuracy, Diversity, Complexity
    \item LoRA es la mejor técnica de SFT en la mayoría de los escenarios; DPO es el algoritmo de alineación más práctico
    \item Model Merging es una herramienta poderosa para combinar las ventajas de varios modelos sin costo adicional
    \item La correlación entre la preferencia humana y los benchmarks automatizados es baja, por lo que es necesario combinar varios métodos de evaluación
    \item Test-Time Compute Scaling intercambia tiempo de inferencia por calidad de salida, y es la tendencia más importante en este momento
\end{itemize}
\end{importantbox}

\subsection{Tabla de decisiones de ingeniería para elegir método}

\begin{table}[H]
\centering
\begin{tabular}{p{0.2\textwidth} p{0.24\textwidth} p{0.24\textwidth}}
\toprule
Escenario objetivo & Método prioritario & Criterio de decisión clave \\
\midrule
Validar primero rápidamente la necesidad & In-context learning / RAG & No pasar demasiado pronto al fine-tuning; primero ver si el base model ya es suficientemente bueno \\
Se necesita un formato de salida o un tono estables & SFT (generalmente con LoRA) & Si la calidad de las muestras es suficientemente alta y si el estilo de salida puede demostrarse de manera explícita \\
Se necesita que las respuestas se ajusten más a la preferencia humana & DPO / Preference Alignment & Si se cuenta con datos de preferencia chosen vs rejected, y si la evaluación cubre la preferencia real de los usuarios \\
Se necesita combinar las especialidades de varios modelos & Model Merging & Si las capacidades de los modelos son complementarias, y si existe riesgo de degradación tras la combinación \\
Aún queda presupuesto en la etapa de inferencia & Test-Time Compute Scaling & Si la latencia y el costo permiten intercambiar más muestreo por mayor calidad \\
\bottomrule
\end{tabular}
\caption{Orden de elección de los principales métodos de Post-Training}
\end{table}

\subsection{Orden de implementación para quien lo lleva a la práctica}

Si se convierte esta clase en recomendaciones de ejecución de ingeniería, el orden más seguro suele ser:

\begin{enumerate}
    \item Primero usar prompting / RAG, de costo mínimo, para delimitar el problema
    \item Después usar un conjunto de datos pequeño y de alta calidad para hacer SFT y resolver los problemas de estilo y formato
    \item Si se necesita mejorar aún más la usabilidad, introducir datos de preferencia para hacer DPO
    \item Por último, según el costo, la latencia y el resultado, elegir model merging o test-time scaling como medio de mejora
\end{enumerate}

\subsection{Lecturas adicionales}

\begin{itemize}
    \item \textbf{LLM Engineer's Handbook}: libro de referencia sistemático escrito por el ponente
    \item \textbf{Open Perfect Blend} (Hugging Face): esquema de mezcla de datos de SFT de propósito general
    \item \textbf{mergekit}: herramienta de código abierto para la combinación de modelos
    \item \textbf{TRL} (Hugging Face): la biblioteca de herramientas de Post-Training más completa
    \item Rafailov et al., \textit{Direct Preference Optimization} (NeurIPS 2023): artículo original de DPO
    \item Hu et al., \textit{LoRA: Low-Rank Adaptation of Large Language Models} (ICLR 2022): artículo original de LoRA
    \item Snell et al., \textit{Scaling LLM Test-Time Compute} (2024): fundamento teórico de Test-Time Compute
    \item Chatbot Arena (\url{https://arena.lmsys.org}): clasificación de modelos basada en la preferencia humana
\end{itemize}

%% --- Fin del contenido --- %%

\end{document}
